\chapter{Marco teórico o Estado del arte}
\label{cap:estado}

Elementos conceptuales que sirven de base al trabajo y estudios previos
relacionados \parencite{smith1980, chapman1981}.

\section{Ejemplo de figura}
Las figuras se numeran con el capítulo (figura~\ref{fig:ejemplo}) y llevan
siempre la fuente debajo.

\begin{figure}[htbp]
  \centering
  \includegraphics[width=0.6\textwidth]{example-image}
  \caption{Modelo de ejemplo}
  \label{fig:ejemplo}
  \fuente{elaboración propia.}
\end{figure}

\section{Ejemplo de tabla}
\begin{table}[htbp]
  \centering
  \caption{Comparativa de motores gráficos}
  \label{tab:motores}
  \begin{tabular}{lcc}
    \toprule
    Motor   & Lenguaje & Licencia \\
    \midrule
    Unity   & C\#      & Propietaria \\
    Godot   & GDScript & MIT \\
    Three.js& JavaScript & MIT \\
    \bottomrule
  \end{tabular}
  \fuente{elaboración propia.}
\end{table}
